% ============================================================
% preambulo-revista-estetica.tex
% ------------------------------------------------------------
% CAPA REEMPLAZABLE DE LAS REVISTAS (ETAPA A1, 07/11/2026). ES LA
% QUE UNA REVISTA CAMBIA POR LA SUYA CUANDO QUIERE SU IDENTIDAD
% VISUAL. NO DEBE DEFINIR NINGUNO DE LOS ENTORNOS NI DE LAS MACROS
% DE LOS CONTRATOS (preambulo-contrato-comun.tex,
% preambulo-revista-contrato.tex): ESOS SE REDEFINEN, SI HACE FALTA,
% CON \renewcommand.
% ------------------------------------------------------------
% SALE DEL PREÁMBULO QUE m_XML.ObtenerPreambuloEmbebido ESCRIBÍA
% DENTRO DE GAMBAS PARA EL CAMINO VIEJO. SE QUITÓ TODO LO QUE AHORA
% DAN LOS CONTRATOS (babel, csquotes, biblatex, siunitx, CUADROS,
% RECUADROS, CONVERSACIÓN, VERSO, CÓDIGO, CORTE Y NOTA EN TÍTULOS,
% TACHADO). DIFERENCIAS VISIBLES CON EL CAMINO VIEJO, POR TOMAR LA
% DEFINICIÓN COMÚN CON LIBROS:
%   - gbverso: SIN \leftmargini=14pt NI \normalsize.
% LAS FUENTES SON LAS DE LIBROS (preambulo-estetica.tex): LA FORMA DE
% CARGAR IBM Plex Mono DEL PREÁMBULO VIEJO NOMBRABA UN ARCHIVO
% IBMPlexMono.otf QUE NO EXISTE.
% ------------------------------------------------------------
% LAS PRESTACIONES, COMO EN LIBROS, SON BANDERAS DE
% preambulo-numero.tex: \ifgbshowframe (LA CAJA A LA VISTA) Y
% \ifgbhomeoarchy (LAS MARCAS DE impnattypo).
% LA PRIMERA PÁGINA (A1.1b, CONTRATO DE REVISTA 2): LA MECÁNICA ES
% DEL CONTRATO; ACÁ, SU ASPECTO (3b) Y LA CABECERA firstpage.
% ============================================================

% ============================================================
% 1. TIPOGRAFÍA (LAS MISMAS FAMILIAS QUE LIBROS)
% ============================================================
\usepackage{fontspec}
\usepackage[final,babel=true]{microtype}
\microtypecontext{spacing=nonfrench}

\def\gbfuentetexto{Libertinus Serif}
\def\gbfuentetitulos{IBM Plex Sans Condensed}
\def\gbfuentemono{IBM Plex Mono}

\setmainfont{\gbfuentetexto}[
  Numbers={OldStyle,Proportional},
  Ligatures={TeX,Common,Discretionary},Scale=1.18]
\setsansfont{\gbfuentetitulos}[
  Scale=MatchLowercase,
  Ligatures=TeX,
  Extension=.otf,
  UprightFont=*-Regular,
  ItalicFont=*-Italic,
  BoldFont=*-SemiBold,
  BoldItalicFont=*-SemiBoldItalic]
\setmonofont{\gbfuentemono}[
  Scale=0.91,
  Extension=.otf,
  UprightFont=*-Regular,
  ItalicFont = IBMPlexMono-Italic.otf,
  BoldFont = IBMPlexMono-Bold.otf,
  BoldItalicFont = IBMPlexMono-BoldItalic.otf]

\usepackage{unicode-math}
\setmathfont{Libertinus Math}[Scale=MatchLowercase]

\renewcommand{\normalsize}{\fontsize{10pt}{13pt}\selectfont}

% ============================================================
% 2. PÁGINA: A4, TEXTO A LA IZQUIERDA Y UNA COLUMNA AL MARGEN
% ------------------------------------------------------------
% headheight 29,5 mm: LA CABECERA DE firstpage (LOGO DE 20 mm, 3 mm,
% FRANJA DE 6 mm) MIDE 82,91 pt Y CON 29 mm (82,51 pt) fancyhdr
% AVISABA. EL MEDIO MILÍMETRO SALE DE top (13 → 12,5): LA CAJA DE
% TEXTO EMPIEZA DONDE EMPEZABA (includehead).
% ============================================================
\ifgbshowframe
  \usepackage[paperwidth=210mm,paperheight=297mm,
    top=12.5mm,bottom=23mm,left=13mm,right=87mm,
    footskip=12mm,headheight=29.5mm,headsep=10mm,
    marginparsep=6mm,marginparwidth=68mm,
    includehead,includefoot,showframe]{geometry}
\else
  \usepackage[paperwidth=210mm,paperheight=297mm,
    top=12.5mm,bottom=23mm,left=13mm,right=87mm,
    footskip=12mm,headheight=29.5mm,headsep=10mm,
    marginparsep=6mm,marginparwidth=68mm,
    includehead,includefoot]{geometry}
\fi
\setcounter{secnumdepth}{0}


% ============================================================
% 3. COLOR
% ------------------------------------------------------------
% gbcolorbase ES EL COLOR DE LOS RECUADROS DEL CONTRATO COMÚN
% ============================================================
\definecolor{azulrevista}{HTML}{2d5a8e}
\colorlet{gbcolorbase}{azulrevista}

% ============================================================
% 3b. PRIMERA PÁGINA: LOS GANCHOS DEL CONTRATO DE REVISTA 2 (4.2)
% ------------------------------------------------------------
% EL ASPECTO DEL DISEÑO VIEJO: ROTULOS EN AZUL CLARO, AUTORES EN AZUL,
% RESÚMENES EN UN RECUADRO SIN FILETE COLOR TÉ CON LECHE
% ============================================================
\renewcommand{\gbestilorotulo}{\bfseries\color{azulrevista!70}}
\renewcommand{\gbestiloautor}{\bfseries\color{azulrevista}}
% LAS DIRECCIONES DE LA COLUMNA, EN SU LETRA: \UrlFont (7) ES LA DEL
% CUERPO
\renewcommand{\gbestilocolumna}{\sffamily\small\RaggedRight\urlstyle{same}}
\tcbset{gbresumen/.style={colback=brown!8!white,colframe=brown!8!white,
  boxrule=0pt,arc=2pt,left=4pt,right=4pt,top=3pt,bottom=3pt,
  before skip=4pt,after skip=4pt}}

% ============================================================
% 4. FIGURAS, NOTAS Y LISTAS
% ============================================================
\usepackage[labelfont=bf,font=small,labelsep=period,format=plain]{caption}

\usepackage[bottom,stable,hang]{footmisc}
\usepackage{scrextend}
\setlength{\footnotesep}{10pt}
\renewcommand{\footnoterule}{%
  \kern -3pt%
  \hrule height 0.5pt width 0.4\columnwidth%
  \kern 6pt}
\renewcommand*{\thefootnote}{\scriptsize\sf{[\arabic{footnote}]}}
\makeatletter
\patchcmd{\@footnotetext}{\footnotesize}{\small}{}{}
\makeatother
\newcommand*\footnotemarkspace{0em}
\deffootnote{\footnotemarkspace}
  {\parindent}
  {\makebox[\footnotemarkspace][r]{\llap{\thefootnotemark\quad}}}

\raggedbottom
\clubpenalty=10000
\widowpenalty=10000

\usepackage{enumitem}
\setlist{nosep,topsep=4pt}

\renewenvironment{quote}
  {\normalsize\list{}{\sf\leftmargin=14pt\rightmargin=0pt}
   \item\relax}{\endlist}

\ifgbhomeoarchy
  \usepackage[hyphenation,homeoarchy,homeoarchywordcolor=orange,
    homeoarchycharcolor=orange,draft]{impnattypo}
\fi

% ============================================================
% 5. TÍTULOS DE SECCIÓN
% ============================================================
\usepackage{titlesec}
\titleformat{\section}[hang]
  {\sf\bfseries\Large\color{azulrevista}\raggedright}
  {\thesection}{.5em}{}[]
\titlespacing{\section}
  {\parindent}{18pt plus 0.5pt minus 0.5pt}{6.75pt}
\titleformat{\subsection}[hang]
  {\sf\large\color{azulrevista}\raggedright}
  {\thesubsection}{.5em}{}[]
\titlespacing{\subsection}
  {\parindent}{18pt plus 0.5pt minus 0.5pt}{6.75pt}
\titleformat{\subsubsection}[hang]
  {\rm\bfseries\normalsize\color{azulrevista}\raggedright}
  {\thesubsubsection}{.5em}{}[]
\titlespacing{\subsubsection}
  {\parindent}{18pt plus 0.5pt minus 0.5pt}{6.75pt}
\titleformat{\paragraph}[runin]
  {\bfseries\color{azulrevista}}{}{0em}{}[\mbox{ --- }]
\titlespacing{\paragraph}
  {0pt}{6.75pt plus 0.5pt minus 0.5pt}{0pt}
\titleformat{\subparagraph}[hang]
  {\sf\bfseries\large\color{azulrevista}\centering}
  {\thesubparagraph}{.5em}{}[]
\titlespacing{\subparagraph}
  {\parindent}{18pt plus 0.5pt minus 0.5pt}{6.75pt}

% ============================================================
% 6. CABECERA Y PIE
% ------------------------------------------------------------
% firstpage: LA DE LA PRIMERA PÁGINA, CON EL LOGO DE LA REVISTA
% (media/logo-revista.png DEL PROYECTO) Y LA FRANJA DEL TIPO DE
% ARTÍCULO. LAS DEMÁS: FOLIO Y DATOS DE LA REVISTA ARRIBA.
% ============================================================
\usepackage{fancyhdr}
\setlength{\headwidth}{\textwidth}
\addtolength{\headwidth}{\marginparsep}
\addtolength{\headwidth}{\marginparwidth}
\fancypagestyle{firstpage}{%
  \fancyhf{}%
  \fancyhead[L]{%
    \begin{minipage}[t]{\headwidth}%
      \includegraphics[width=184mm,height=20mm,keepaspectratio=false]{media/logo-revista.png}\\[3mm]%
      {\setlength{\fboxsep}{0pt}%
      \colorbox{azulrevista}{%
        \parbox[c][6mm][c]{\headwidth}{%
          \hspace{12pt}\sffamily\bfseries\color{white}\articulotipo%
        }%
      }}%
    \end{minipage}%
  }%
  \renewcommand{\headrulewidth}{0pt}%
  \fancyfoot[L]{%
    \begin{minipage}[t]{\headwidth}%
      \textcolor{azulrevista}{\rule{\headwidth}{1pt}}\\[1pt]%
      \sffamily\thepage\hfill%
      \small\articulorevista{} / e-ISSN: \texttt{\articuloissne}{} / Vol.~\articulovolumen{},~n.\textsuperscript{o}~\articulonumero{} / \articulofecha%
    \end{minipage}%
  }%
  \renewcommand{\footrulewidth}{0pt}%
}
\pagestyle{fancy}
\fancyhf{}
% LAS INTERIORES: FOLIO Y DATOS DE LA REVISTA; DEBAJO, LOS AUTORES
% (REDALYC 4.5: EL MEMBRETE DE LAS INTERIORES LLEVA LOS AUTORES)
\fancyhead[L]{%
  \raisebox{\dimexpr\headheight-7mm\relax}[0pt][0pt]{%
    \begin{minipage}[b]{\headwidth}%
      \sffamily\thepage\hfill%
      \small\articulorevista{} / e-ISSN: \texttt{\articuloissne}{} / Vol.~\articulovolumen{},~n.\textsuperscript{o}~\articulonumero{} / \articulofecha\\[1pt]%
      \mbox{}\hfill{\small\itshape\articuloautorescorto}\\[1pt]%
      \textcolor{azulrevista}{\rule{\headwidth}{1pt}}%
    \end{minipage}%
  }%
}
\renewcommand{\headrulewidth}{0pt}

% ============================================================
% 7. CORTE DE DIRECCIONES ELECTRÓNICAS
% ------------------------------------------------------------
% LA MISMA REGLA QUE LIBROS (preambulo-estetica.tex): SE PERMITE EL
% CORTE EN CUALQUIER CARÁCTER. url SE CARGA ANTES QUE hyperref:
% \Urlmuskip NO EXISTE HASTA QUE EL PAQUETE ESTÁ
% ============================================================
\usepackage{url}
\Urlmuskip = 0mu plus 1mu
\def\UrlBreaks{\do\a\do\b\do\c\do\d\do\e\do\f\do\g\do\h\do\i\do\j\do\k\do\l\do\m\do\n\do\o\do\p\do\q\do\r\do\s\do\t\do\u\do\v\do\w\do\x\do\y\do\z\do\A\do\B\do\C\do\D\do\E\do\F\do\G\do\H\do\I\do\J\do\K\do\L\do\M\do\N\do\O\do\P\do\Q\do\R\do\S\do\T\do\U\do\V\do\W\do\X\do\Y\do\Z\do0\do1\do2\do3\do4\do5\do6\do7\do8\do9\do=\do.\do:\do\%\do?\do_\do-\do+\do/\do\#\do~}
\def\UrlFont{\rmfamily}

% ============================================================
% 8. ENLACES (ÚLTIMO)
% ------------------------------------------------------------
% hyperref VA AL FINAL, COMO EN LIBROS (preambulo-hyperref.tex). LOS
% DATOS DEL PDF (TÍTULO, AUTORES, DERECHOS, LICENCIA) LOS FIJA EL
% ARTÍCULO CON \hypersetup AL EMPEZAR (jats-a-latex-articulo.xsl)
% ============================================================
\usepackage[
  unicode=true,
  pdfencoding=auto,
  psdextra,
  colorlinks=true,
  linkcolor=azulrevista,
  citecolor=azulrevista,
  urlcolor=azulrevista,
  filecolor=azulrevista,
  breaklinks=true,
  bookmarks=true,
  bookmarksopen=true,
  bookmarksnumbered=true
]{hyperref}
% SIN hyperxmp: EL ARTÍCULO SE COMPILA SIEMPRE CON \DocumentMetadata
% (PDF/A-2u), QUE ESCRIBE EL XMP Y DESACTIVA hyperxmp (MEDIDO: «The
% loading of package hyperxmp is disabled»)
